% Derivation of the Doppler effect on a Gaussian wave packet in the satellite frame
% Based on handwritten notes (data/lorentz.jpg)

\documentclass[12pt,a4paper]{article}
\usepackage[utf8]{inputenc}
\usepackage[T1]{fontenc}
\usepackage[french]{babel}
\usepackage{amsmath, amssymb}
\usepackage{physics}
\usepackage{geometry}
\geometry{margin=2.5cm}

\begin{document}

\section*{Effet Doppler relativiste sur un paquet d'ondes quantique}

\subsection*{Paquet d'ondes gaussien — référentiel terrestre}

On modélise l'impulsion quantique (photon) par un paquet d'ondes gaussien en temps retardé
$\tau = t - x/c$\,:

\begin{equation}
    \Psi(\omega_0;\, \tau)
    = \frac{1}{\sqrt{\pi}\,\tau_c}
      \, e^{-\tau^2/(2\tau_c^2)}
      \, e^{i\omega_0 \tau}
    \label{eq:psi_earth}
\end{equation}

où $\omega_0$ est la fréquence centrale et $\tau_c$ est la durée de cohérence de l'impulsion.
Ce paquet est normalisé :

\begin{equation}
    \int_{-\infty}^{+\infty} \abs{\Psi(\omega_0;\, \tau)}^2 \, d\tau = 1
\end{equation}

\subsection*{Passage dans le référentiel du satellite}

Soit $\beta = v_\mathrm{sat}/c$ le rapport de la vitesse du satellite à la vitesse de la lumière,
et $\gamma = 1/\sqrt{1-\beta^2}$ le facteur de Lorentz. La transformation de Lorentz (boost
selon l'axe de propagation $x$) s'écrit :

\begin{equation}
    \begin{cases}
        t' = \gamma\!\left(t - \dfrac{\beta\, x}{c}\right) \\[8pt]
        x' = \gamma\!\left(x - \beta c\, t\right)
    \end{cases}
    \qquad \Longleftrightarrow \qquad
    \begin{cases}
        t  = \gamma\!\left(t' + \dfrac{\beta\, x'}{c}\right) \\[8pt]
        x  = \gamma\!\left(x' + \beta c\, t'\right)
    \end{cases}
    \label{eq:lorentz}
\end{equation}

\subsection*{Transformation du temps retardé}

Le temps retardé dans le référentiel du satellite est $\tau' = t' - x'/c$. En appliquant
\eqref{eq:lorentz} :

\begin{align}
    \tau' &= t' - \frac{x'}{c}
           = \gamma\!\left(t - \frac{\beta x}{c}\right)
             - \frac{\gamma}{c}\!\left(x - \beta c\, t\right) \notag \\[4pt]
          &= \gamma t(1 + \beta) - \frac{\gamma x}{c}(1 + \beta)
           = \gamma(1+\beta)\!\left(t - \frac{x}{c}\right)
\end{align}

On obtient donc la relation fondamentale :

\begin{equation}
    \boxed{
        \tau' = \gamma(1+\beta)\,\tau
              = \sqrt{\frac{1+\beta}{1-\beta}}\;\tau
    }
    \label{eq:tau_boost}
\end{equation}

ce qui s'inverse en :

\begin{equation}
    \tau = \gamma(1-\beta)\,\tau'
         = \sqrt{\frac{1-\beta}{1+\beta}}\;\tau'
    \label{eq:tau_inverse}
\end{equation}

\subsection*{Paquet d'ondes dans le référentiel du satellite}

On substitue \eqref{eq:tau_inverse} dans \eqref{eq:psi_earth}. Le changement de variable
$\tau \to \tau'$ introduit un jacobien $d\tau/d\tau' = \sqrt{(1-\beta)/(1+\beta)}$, ce qui
modifie la normalisation. Pour conserver $\int|\Psi'|^2 d\tau' = 1$, il faut multiplier par
$\sqrt{d\tau/d\tau'} = \left(\frac{1-\beta}{1+\beta}\right)^{1/4}$.

On obtient le paquet d'ondes dans le référentiel du satellite :

\begin{equation}
    \boxed{
    \Psi'(\omega_0;\, \tau')
    = \frac{1}{\sqrt{\pi}\,\tau_c}
      \left(\frac{1-\beta}{1+\beta}\right)^{\!1/4}
      \exp\!\left[-\frac{1-\beta}{1+\beta}\,\frac{\tau'^2}{2\tau_c^2}\right]
      \exp\!\left[i\sqrt{\frac{1-\beta}{1+\beta}}\;\omega_0\,\tau'\right]
    }
    \label{eq:psi_sat}
\end{equation}

\subsection*{Interprétation}

L'expression \eqref{eq:psi_sat} est encore un paquet d'ondes gaussien, avec les paramètres
décalés par effet Doppler relativiste. En définissant le facteur Doppler
$D = \sqrt{(1+\beta)/(1-\beta)}$, on identifie :

\begin{align}
    \omega_0' &= \frac{\omega_0}{D} = \omega_0\sqrt{\frac{1-\beta}{1+\beta}}
    && \text{(décalage vers le rouge pour un satellite qui s'éloigne)} \\[6pt]
    \tau_c'   &= D\,\tau_c = \tau_c\sqrt{\frac{1+\beta}{1-\beta}}
    && \text{(dilatation de la durée d'impulsion)}
\end{align}

de sorte que \eqref{eq:psi_sat} se réécrit sous la forme standard :

\begin{equation}
    \Psi'(\omega_0';\, \tau')
    = \frac{1}{\sqrt{\pi}\,\tau_c'}
      \, e^{-\tau'^2/(2\tau_c'^2)}
      \, e^{i\omega_0'\,\tau'}
\end{equation}

La structure gaussienne du paquet est donc préservée par la transformation de Lorentz,
avec une fréquence et une durée de cohérence toutes deux multipliées par le facteur Doppler.

\end{document}
